\documentclass[12pt]{article}
\usepackage{hyperref}

\usepackage{graphicx}

\usepackage{mathtools}
\DeclarePairedDelimiter\ceil{\lceil}{\rceil}
\DeclarePairedDelimiter\floor{\lfloor}{\rfloor}

\usepackage{cancel}

\usepackage{amsmath}
\usepackage{amsfonts}
\usepackage{amssymb}

\usepackage{mathtools}
\usepackage{xcolor}
\usepackage{listings}
\definecolor{darkgreen}{rgb}{0.0, 0.2, 0.13}
\lstset{language=C++,
                basicstyle=\ttfamily,
                keywordstyle=\color{blue}\ttfamily,
                stringstyle=\color{red}\ttfamily,
                commentstyle=\color{darkgreen}\ttfamily,
                morecomment=[l][\color{magenta}]{\#}
}

\usepackage{breqn}
\usepackage[]{algorithm2e}
\usepackage{physics}

\newcommand\fnurl[2]{%
  \href{#2}{#1}\footnote{\url{#2}}%
}

\newcommand{\Four}{\mathcal{F}}
\renewcommand{\(}{\left(}
\renewcommand{\)}{\right)}
%% \renewcommand{\{}{\left{}
%% \renewcommand{\}}{\right}}


\title{Spectral Operation design doc for rocFFT}
\author{Malcolm Roberts and Steve Leung}



\def\no#1{\nobreak{#1}} % stop equation breaking in dmath

\begin{document}
\maketitle

\section{Use cases}

\subsection{Shuffled output}

If we do a pointwise operation that depends on the index, how do
we provide this information to the user-specified operator?

\subsection{Convolution/correlation}

We can have an arbitrary number of inputs.  Each input can have its
own format (length, stride, precision).  Typically, inputs will have
the same type (ie real vs complex), but it would be perfectly
reasonable to convolve real data with complex data, or complex data
with Hermitian-symmetric complex data.  Real and Hermitian-symmetric
complex wouldn't really make sense, as they are on opposite sides of
the Fourier transform.  In general, we expect that all inputs and
outputs be in the same space, ie either all in the time/space domain
or all in the frequency domain.

The convolution of two inputs $\{f\}_{n=0}^{N-1}$ and  $\{g\}_{n=0}^{N-1}$
is denoted $f\star{}g$ and is given by
\begin{dmath}
  f\star{}g_\ell
  = \sum_{m=0}^{n}f_m g_{\ell-m}
  = \sum_{m_0=0}^{N-1}\sum_{m_1=0}^{N-1}f_{m_0} g_{m_1} \delta_{\ell, m_0 + m_1}
\end{dmath}
where
\begin{dmath}
  \delta_{a,b} = \begin{cases}
  1  & a=b  \\
  0 & a \neq b
\end{cases}
\end{dmath}
is the Kronecker delta.  This can be extended from a binary convolution
to a general $n$-ary conovlution as
\begin{dmath}
  \star{}\(f_0, \dots, f_{n-1} \)_\ell
  = \sum_{m_0, \dots, m_{n-1}} f_{m_0} \dots f_{m_{n-1}}  \delta_{\ell, \sum m_i} 
\end{dmath}

Hermitian input needs a different algorithm depending on the order of
the operator.

So, for two given inputs $f$ and $g$ of length $N$ and $M$
\href{https://docs.scipy.org/doc/scipy/reference/generated/scipy.signal.fftconvolve.html#scipy.signal.fftconvolve}{python}
uses
\begin{itemize}
\item \texttt{full}: $\max(M,N)$
\item \texttt{valid}: $M+N-1$
\item \texttt{same}:  $\max(M, N) - \min(M, N) + 1$.
\end{itemize}

(Hockney Free Space Convolution?)

\subsection{Shift}

Multiply by roots of unity.  Vanderlinde asked for a batched 1D FFT
with per-batch offset, ie
\begin{dmath}
  u_{n,b} \leftarrow \omega^{\alpha_b + \beta n} u_{n,b}
\end{dmath}
where $b$ is the batch index and $n$ is the input index.  NB: this is
annoying for real-to-complex transforms.


\subsection{Poisson solver}

Division by $k^2$.

\subsection{filter}

high/low-pass filter.  NB: this can be further optimised by not doing
inverse transforms on known-zero data, but, well.

\subsection{Vanderlinde}

wants an affine shift, and a norm per batch.  The second is admitted
as being kind of unrealistic.


\subsection{applications}

upsampling: looks like a F, zero-padding, Finv?
(\href{https://stackoverflow.com/questions/5156690/help-with-resampling-upsampling}{stack
  overflow link})

SAR: synthetic aperture radar

M\_TIP SPI reconstruction algorithm

non-uniform FFT  (Flatiron institute)

Window'd transforms

image compression

How about finding the frequency offset?

Shuffled output

Fractional-dimensional transforms?

Resample: pad with zeros in the spectral domain, then transform back
(longer output?)  This seems like a hard case to manage.

\section{Other libraries}

vkFFT

cuFFTDx

Spiral

FFTW++

cutensor: good way to get the data in there?
\url{https://docs.nvidia.com/cuda/cutensor/getting_started.html}

\section{API proposals}

There are two basic cases:
\begin{enumerate}
\item Part-of-transform
\item Round-trip
\end{enumerate}

\subsection{Indexing function}

The pointwise operation may depend on the logical index, which isn't
necessarily going to be the same as the array index.  For example, 1D
transforms which use multiple kernels will have effectively transposed
the data; the map from array index to logical index depends on the FFT
plan.  For index-dependant functions, we will need the index map in
order to correctly apply the pointwise function.

\subsection{Part-of-transform}

For differentiation, shift, result scaling, and amplitude filters.

The functionality here is similar to use cases for callback API.
However, the callback functionality doesn't know about the
array-index-to-logical-index map.  Also, the performance of callbacks
have been very problematic on the GPU, so moving to another API where
we can RTC user-specified functions into the FFT kernel will allow for
significantly better performance.  As with the callback API, we can specify
a load and a store function.  The functions have signature
\begin{lstlisting}
Tval f(Tval u, size_t idx_x, size_t idx_y, size_t idx_z)
\end{lstlisting}
That is, they take exactly one value (real or complex) and a
multi-index, and return exactly one value of the same type.

For example, a Poisson solve is provided by
\begin{lstlisting}
Tval f(Tval u, size_t idx_x, size_t idx_y, size_t idx_z)
{
  const double k0 = 2.0; // User-specified
  const double k2 = k0 * k0
         * (idx_x * idx_x +  idx_y * idx_y +  idx_z * idx_z);
  if(k2 == 0)
     return u;
  return u / k2;
}
\end{lstlisting}

A low-pass filter is provided by
\begin{lstlisting}
Tval f(Tval u, size_t idx_x, size_t idx_y, size_t idx_z)
{
  const double k0 = 2.0; // User-specified
  const double k2 = k0 * k0
         * (idx_x * idx_x +  idx_y * idx_y +  idx_z * idx_z);
  if(k2 > cutoff)
     return 0;
  return u;
}
\end{lstlisting}

A shift operator is provided by
\begin{lstlisting}
Tval f(Tval u, size_t idx_x, size_t idx_y, size_t idx_z)
{
  const Tval omega = exp(2 * pi / N); // User-specified
  return u * pow(omega, a + b * idx_x);
}
\end{lstlisting}

We will have to provide a method of getting the complex type from the
real type and vice-versa.

\subsection{Round-trip}

Convolution, Poisson solve, high/low-pass filters, 

Can be combined with pre/post functions, eg: windowed transform.

We add the information to the \texttt{rocfft_plan_description}.  We
need to specify the lengths and strides separately for each input
buffer.  We can have an arbitrary number of inputs and outputs.  Outputs
can be mixed real/complex or complex/Hermitian-complex.

We need to specify the order of the operation: Poisson solves are
order-1, binary convolutions are order-2, ternary convolutions are
order-3, etc.

We need to specify the pointwise operator.  This has to match the
order of the operation (eg: $-u/k^2$ is order-1, $u\odot v$ is
order-2, etc).  If the operator has higher order than what is
specified, then incorrect results will be returned.

We need to specify the size of the output.  TODO: this isn't super
clear what options there are; cf numpy outputs.  This also has to be
generalized to higher-order transforms, and also made clear for
first-order operations.

TODO: add functions for Navier--Stokes, MHD, Maxwell's equations.
(And divergence-free optimized versions.)  Nonlinear Schroedinger?

Question: if we have complex data, does it matter if we do a forward
or inverse transform?  (is there a complex-conjugation issue?)


\begin{lstlisting}
rocfft_status rocfft_roundtrip_plan_create(...)
\end{lstlisting}

\begin{lstlisting}
rocfft_status rocfft_roundtrip_execute(
  const rocfft_plan     plan,
  void** in_buffers,
  void** out_buffers)
\end{lstlisting}


\subsection{Interface for pointwise function}

\begin{enumerate}
\item Just use callbacks?  Super slow
\item User-provided string?  (Check that it fits the API by doing a
  compilation pre-test in a sandbox environment.)
  \subitem Presets (with enum?)
  \subitem Examples?  With common cases provided?
\item Allow users to select from a list of rocFFT-provided options?
  (enums?)
\item rocRoller?
\end{enumerate}

If we have, for example, a shift operator followed by a filter
operator, who is responsible for ensuring that the new index matches
the filter operator's indexing scheme?

\section{Other infra}

hipGraph?

Multi-gpu?






\end{document}
